% SOURCE docs/release/2026-10-29/source/docs/writing-repair/2026-10-29/src/bf867d2e04e8dbb7.tex CHAR 0
\chapter{Choosing the next experiment: information, decisions and agent limits}
\label{exec13:active-experiments}

An agent that proposes a plausible causal diagram has not yet answered the
bank's next question: which additional observation could change a worthwhile
decision? A diagram may contain a dozen unresolved arrows, although only one
uncertainty affects whether a particular retention offer should be deployed.
Conversely, an apparently stable effect estimate can conceal a decisive
measurement or selection problem. Experiment selection therefore connects
problem formulation, causal discovery and investment analysis. Its purpose is
to spend a limited information budget on distinctions that matter, while
preserving the possibility that none of the proposed experiments is useful.

The proposed role of a language model is to formulate competing mechanisms,
translate a question into an admissible design menu, locate observations that
could distinguish the mechanisms, and explain the resulting numerical
comparison. A statistical component must evaluate the observation laws and
decision losses. A language model's confidence, the number of agreeing agents,
and the novelty of an experiment description are inadequate substitutes for
those calculations. This chapter develops that separation using inspected
primary literature, concrete public software interfaces, and an exact banking
planning example. The example establishes a property of its stipulated law;
it does not measure language-model ability or the value of a bank experiment.

\section{Three different meanings of an experimental action}

A real intervention changes an admissible treatment or environment and obtains
an outcome under that change. A simulator experiment changes an input in a
specified mathematical world. A model or expert query obtains an answer from
a fallible information source. All three can be sequential actions, but their
observation laws have different meanings. A simulator can reveal its own
response to an offer without establishing the response of depositors. An
expert can state that a variable precedes another without supplying a
randomized treatment assignment. Neither operation becomes a bank intervention
because an agent calls it an experiment.

Let $\theta$ index a candidate mechanism, $D$ the information available before
design selection, $d\in\mathcal D(D)$ an admissible acquisition, and $O_d$ its
returned observation. A Bayesian design calculation requires
$p(O_d\mid\theta,d,D)$ and a distribution $\pi(\theta\mid D)$.
For a physical experiment, the former must represent assignment, adherence,
measurement, missingness and the outcome clock. For a simulated experiment,
it is the generator's observation law. For an expert query, it is an answer
error model, potentially conditional on wording, context and common source
material. Estimating a histogram of repeated answers measures a query
distribution; without additional calibration it does not identify the
probability that an asserted causal arrow is correct.

Consider a bank deciding whether to offer a fee waiver to a defined group of
depositors. The target is incremental 90-day stable-balance value, net of the
waiver, contact and funding costs, relative to the existing policy. Define
the population before contact, use only its pre-contact information for
eligibility, and distinguish the offer from acceptance of the offer. A
proposed prospective study would need an assignment rule and a way to
observe both withdrawals and account closure. A historical query about
customers who accepted cannot silently take the place of that study.
Selection, interference across accounts and changes in market rates could
each invalidate the proposed observation law.

The information clock is as consequential as the treatment clock. A design
chosen today may produce a 90-day outcome before a later cohort's deployment
decision; it cannot inform today's irreversible decision before that outcome
exists. Transport from the study cohort to the deployment cohort is then an
additional assumption. Delay and the value of waiting belong in the cost
calculation. Public or generated data can test whether an agent respects these
distinctions now, even when conducting the prospective bank experiment would
be premature.

\section{What the primary literature supports}

\subsection{Learning a strategy does not remove its information requirements}

Lampinen and colleagues study whether agents can acquire active causal
strategies from passive demonstrations. The positive finding concerns learning
how to investigate from examples of competent investigation. At evaluation,
the learner can act and obtain the consequences needed by the task. Passive
training is therefore compatible with active information acquisition; it is
not evidence that an observationally unidentified effect becomes identified
without new assumptions or observations. The paper's finite graphs and
constrained action and observation spaces also limit transfer to a bank with
latent customer states, delayed outcomes and restricted actions
\citep{v12_la_2305_16183}.

This distinction suggests a useful proposal component. Demonstration traces
can teach an agent to ask for the outcome horizon, request a missing
assignment variable, or select a discriminating synthetic intervention. Their
training value should be evaluated separately from the information supplied
by the eventual environment. A comparison should hold that environment and
its acquisition budget fixed, and include a numerical or rule-based design
policy that receives the same history. Otherwise a more informative interface
can be mistaken for a better reasoner.

Verma and colleagues provide a different example: sequential causal discovery
with a language model treated as a noisy expert. Their method combines batches
of observational data with background information obtained through a fixed
budget of model queries, using partial ancestral graphs to retain relevant
latent-confounding ambiguity. The inspected evaluation compares query
allocation with random edge selection and reports improvements in graph
metrics. This is evidence about the tested graph-learning setup, not a
physical intervention or a measured gain from a lending decision. Sequential
arrival of data batches also does not by itself supply the treatment-history
assumptions required by longitudinal causal inference
\citep{v12_la_2506_16234}.

The paper identifies open questions that directly affect deployment: the
uncertainty over candidate graphs, calibration of language-model answers,
more complex latent confounding, the incorporation of interventions and
sequential theoretical analysis. These limits argue for retaining ambiguous
graphs and for pricing questions according to their possible effect on the
decision. Repeated model answers may share training information and reasoning
errors. An empirical answer frequency is useful operational evidence, but
the product should not multiply it into independent scientific confirmation.

\subsection{Model discovery and numerical design: gains with important exceptions}

The inspected version of Model Discovery Agent combines model proposals with
numerical model comparison and experimental design. Candidate mechanistic
models are fit to observations, and a design component seeks experiments whose
predictions distinguish the current candidates. The paper discusses Bayesian
model evidence, approximations including BIC, sequential Monte Carlo, and
disagreement-based surrogates for expected information gain. Depending on the
design space, its implementation uses numerical optimization or evaluates a
finite menu. This combination offers a concrete reason to use a language
model: it can propose representations and mechanisms that a fixed estimator
does not enumerate, while numerical components fit and compare the proposals
\citep{v12_la_2608_09696}.

The favorable evidence needs its stated denominator and comparison budget.
The paper reports experiments in chemistry, glucose dynamics and neuronal
systems; its chemistry aggregate uses 115 matched evaluable world--seed cells
from 171 attempted cells over 54 worlds. This is a conditional comparison over
evaluable cells, not a claim that all attempts succeeded. The paper also
explicitly distinguishes sample efficiency from a compute-matched comparison.
Extra model proposals, optimization, fitting and retries may consume substantial
resources even when fewer environment queries are used. A bank investment
argument must record both kinds of expenditure.

Adverse findings sharpen the design recommendation. In the paper's
source-location experiment, a greedy joint-variance heuristic performs worse
than random selection, and the model-information criterion does not show an
advantage over random design in that domain. A numerical-looking acquisition
score is therefore not a guarantee of improved information gathering.
The observation model, candidate pool, approximation error and accessible
design menu all affect whether the score is useful. Random and fixed designs
remain substantive baselines; they are not merely weak comparators to be
discarded once an agent is available.

A further limitation is support. Normalized weights over the proposed models
sum to one even when every model omits the mechanism that generated the data.
Pruning can make apparent confidence increase while removing a useful
alternative. For the bank, candidate mechanisms should include policy
selection, timing error, measurement changes and null effects where relevant.
An agent should be able to propose a new mechanism after a failed predictive
check. Such revision needs an evaluation rule that prevents unlimited
post hoc repairs from making every hypothesis unfalsifiable.

\subsection{Scientific judgment, research state and the missing validation step}

DISCERN tests a complementary problem: whether agents carry evidence
limitations through scientific reasoning. Its 203 canonical life-science tasks
comprise 79 data-integrity tasks, 102 analysis-verification tasks and 22
hypothesis-generation dossiers, across eight tracks and eight evaluated models.
The three levels are independently instantiated and scored. Their outputs
are not passed through an evaluated end-to-end research pipeline. For a
counterfactual dossier, eligibility depends on eliciting a model prior that
conflicts with the supplied evidence
\citep{exec13_active_2609_33357}.

The reported perfect-score rates are 60.8\%, 34.2\% and 0.6\% at the three
levels under their respective rubrics. These are rubric-specific evaluation
percentages, not probabilities of discovering a true mechanism or choosing
a profitable experiment. The grading uses human-designed criteria,
deterministic checks with language-model verification and adjudication for
the first two levels, and rubric-based judgment for hypotheses. Three draws
of a judging procedure do not amount to three independent institutions
validating an experimental result. The authors identify task dependence,
possible familiarity with public data and residual grading uncertainty, and
do not establish whether the proposed hypotheses lead to useful follow-up
experiments.

For banking, the relevant failure mechanism is precise. An agent can correctly
say that offer acceptance is selected, yet subsequently recommend deploying
the offer as though the observed acceptance--balance association were a causal
effect. Merely counting whether the word ``confounding'' appears in a report
misses that failure. A useful evaluation checks whether the limitation changes
the next action, narrows the claim, or produces abstention. It also needs clean
controls: an agent that rejects every dataset can appear cautious while
destroying valid opportunities. The proposed experiment-selection evaluation
should score both unwarranted deployment and unnecessary rejection.

Discovery Foundation Models proposes an explicit, revisable research state,
with representations, hypotheses, evidence, experiment choices and transferable
research skills as objects that can change. Its Zetema description separates
proposed expectations from external outcomes and calls for attributing a
failure before promoting a lesson to reusable memory. This is a useful
architectural principle: an invalid measurement should revise the observation
model, whereas a failed mechanistic prediction may require a new hypothesis
\citep{exec13_active_2609_15973}.

The report's GALILEO discussion describes a structured therapeutic-discovery
case and explicitly locates specialized human-executed assays inside the
coupled system. It does not establish a general-purpose discovery system or
broad transfer across domains. The underlying biomedical original is outside
the inspected evidence used here, so the account informs the proposed
architecture without independently validating its biological results.
Adopting a revisable state representation is consequently a design choice;
claiming that the proposed bank agent improves scientific discovery would
require a separate controlled evaluation.

\section{A decision criterion with an explicit option to stop}

Let $a\in\mathcal A$ denote a deployment decision and $u(a,\theta)$ its
incremental value, expressed in a common unit and at a specified horizon.
The value of acting with existing information is
\[
 V(D)=\max_{a\in\mathcal A}\mathbb E[u(a,\theta)\mid D].
\]
If acquisition $d$ costs $c(d,D)$, its one-step net value is
\begin{equation}
 G(d\mid D)=\mathbb E_{O_d\mid d,D}
   \left[\max_{a\in\mathcal A}
   \mathbb E[u(a,\theta)\mid D,O_d,d]\right]-V(D)-c(d,D).
 \label{exec13:active-decision-value}
\end{equation}
Acquisition is worthwhile under this model only if its best admissible net
value is positive. Costs can include the experiment itself, failed attempts,
waiting, analyst effort and computation. Combining per-customer value with a
whole-study cost requires an explicit deployment population or amortization;
otherwise the comparison has inconsistent units.

Expected information gain asks a different question. For a finite mechanism
set, write $\mathcal H(\pi)=-\sum_\theta\pi_\theta\log_2\pi_\theta$.
Then
\[
 I(d\mid D)=\mathcal H\{\pi(\cdot\mid D)\}
  -\mathbb E_{O_d\mid d,D}\mathcal H\{\pi(\cdot\mid D,O_d,d)\}.
\]
This criterion can be appropriate when learning the mechanism is itself the
objective. It need not prioritize uncertainty that changes a particular
decision. The target of information gain matters: information about a nuisance
parameter, an effect sign and a full dynamic model need not have the same
economic value. A remedy is to use a decision-relevant target or evaluate
Equation~\eqref{exec13:active-decision-value} directly, while examining the
sensitivity of both to their assumed likelihoods and losses.

\subsection{An exact fee-waiver planning construction}

Suppose two binary variables $H$ and $Z$ index four equally likely planning
worlds. $H$ represents whether the fee-waiver policy is valuable; $Z$ is an
independent nuisance feature of the environment. Relative to abstention,
deploying has value USD~100 when $H=1$ and USD~$-120$ when $H=0$.
Abstention has value zero in either world. These are stipulated, already net
deployment values per normalized deployment decision, not measured customer
effects. Each information acquisition costs USD~5 in the same normalized
units. The calculation assumes that its observation is available before the
deployment decision and that the world remains relevant to that decision
\citep{exec13_active_exact}.

With prior probability $p=\Pr(H=1)=1/2$, expected deployment value is
$100p-120(1-p)=220p-120=-10$. The optimal initial action is abstention.
The threshold for deployment is $p>6/11$. Perfect information about $H$
would have gross value USD~50, since deployment would occur only in the good
half of the worlds. That USD~50 is an upper bound under the stipulated model,
not a forecast of a feasible experiment's return.

Consider two permitted abstract acquisitions. Design $d_Z$ reveals $Z$
perfectly. Design $d_H$ returns a signal $S\in\{0,1\}$ satisfying
\[
 \Pr(S=H\mid H,Z,d_H)=\frac45.
\]
The signal has conditional error probability $1/5$, independent of $Z$.
This accuracy is an assumed experiment likelihood. It is not a measured
accuracy of an LLM, a hypothetical expert panel or an actual bank trial.
Equal acquisition budgets make the comparison transparent: both designs
use one acquisition costing USD~5.

The first design learns one bit about the full world but leaves the posterior
probability of $H$ at $1/2$. The bank still abstains, so net information value
is USD~$-5$. For the second design, each signal has probability $1/2$.
A positive signal gives posterior $4/5$ and expected deployment value
$220(4/5)-120=56$. A negative signal gives posterior $1/5$ and expected
deployment value $-76$, so the bank abstains. Expected terminal value is
$(1/2)56=28$ and net value is USD~23. Full-world information gain is only
$1-h_2(1/5)\simeq0.278$ bits, where
$h_2(p)=-p\log_2p-(1-p)\log_2(1-p)$.

\begin{center}
\begin{tabular}{lrrr}
\toprule
One-acquisition policy & Information (bits) & Gross USD & Net USD\\
\midrule
Maximize information about $(H,Z)$ & 1.000 & 0 & $-5$\\
Maximize deployment value & 0.278 & 28 & 23\\
Uniform random choice of the two & 0.639 & 14 & 9\\
Do not acquire & 0 & 0 & 0\\
\bottomrule
\end{tabular}
\end{center}

The random policy chooses between the two designs with equal probability,
so its expected net value is $(-5+23)/2=9$. Its expected information gain
also averages the two design gains, conditional on knowing the selected
design. An information criterion targeted specifically at $H$ would select
the useful signal in this example. The construction therefore demonstrates
a mismatch between the full-world learning objective and this deployment
objective, rather than a general failure of information-based design.
The values follow from exact finite sums, independently checked by evaluating
the realized payoff of the signal-dependent deployment policy.

\subsection{Unavailable observations and correlated answers}

Now suppose $d_H$ returns no observation with probability $1/4$, independently
of the world, and every attempt still costs USD~5. On failure the posterior
remains the prior and the bank abstains. Gross expected value becomes
$(3/4)28=21$, and all-attempt net value is USD~16. Restricting the comparison
to successful acquisitions would instead report USD~23. Both calculations
can be arithmetically correct, but they answer different questions. A purchase
or investment decision needs the attempted-work denominator. If failure
depends on $H$, the failure event itself has a likelihood and must enter the
posterior; treating it as independent missingness would be another model
error.

Repeated agreement introduces a separate problem. Suppose five expert
answers are perfectly correlated because they all copy one underlying noisy
answer whose accuracy is $4/5$. After five positive answers, the correct
posterior remains $4/5$. Incorrectly multiplying five independent likelihoods
would produce
\[
 \frac{(4/5)^5}{(4/5)^5+(1/5)^5}=\frac{1024}{1025}.
\]
Different prompts or agent identities do not establish independent error
channels. The remedy is to model the shared source, collect evidence from a
different measurement or intervention when feasible, and compare conclusions
under plausible dependence structures. A majority vote cannot create the
missing assignment evidence.

\subsection{A finite stopping rule and the danger of a narrow model pool}

Suppose at most two useful signals are available and their errors are
conditionally independent given $H$. Always purchasing both gives expected
terminal value $148/5=29.6$ and costs USD~10, leaving net value $98/5=19.6$.
This is below the one-query value of USD~23. A conditional calculation verifies
the stopping decision. After a positive first signal, the best possible gross
increment from a second useful query is $16/5=3.2$; after a negative signal
it is zero. Both are below its USD~5 cost. A query about $Z$ adds no deployment
value at either history. Thus, for this menu and horizon of at most two,
the optimal policy purchases one useful query and stops. This result does
not establish that a myopic stopping rule is optimal for arbitrary longer
horizons or complementary experiments.

Finally, suppose a proposed model pool contains only $H=1$. Its normalized
posterior assigns probability one to good outcomes and predicts USD~100 from
deployment. If the omitted bad world is true, deployment instead yields
USD~$-120$. No entropy calculation confined to the single-model pool detects
the omission. The bank therefore needs sensitivity to model support as well
as uncertainty within the pool. Challenging a favorable mechanism with a
null, adverse or selection-based alternative can be more consequential than
refining its parameters. A formal robust recommendation would specify the
larger set of admissible laws and the associated loss criterion; it cannot
be inferred merely from an agent having considered several verbal stories.

\section{Public software: what can be reused and what remains a design}

\subsection{GlucoseBench: a useful acquisition and evaluation boundary}

The inspected GlucoseBench source is a small adaptation of simglucose for
patient-specific intervention forecasting. Its documented target is the same
simulated patient under new interventions, not generalization to unseen
patients or clinical treatment. The maintained interface supplies two passive
training episodes, a menu of available interventions, an \texttt{observe}
operation with a six-episode total budget, and \texttt{freeze} before held-out
evaluation. The inspected evaluator reports failed query cells explicitly;
the summary returns no mean score if any cell in its group failed
\citep{exec13_active_glucosebench}.

These are useful design patterns for the proposed bank library. A separate
evaluator should retain the generator truth, give the learner only the
permitted history and action menu, and prevent held-out outcomes from steering
acquisition. Freezing the acquisition history before scoring makes an
otherwise subtle form of information leakage visible. The in-process object
boundary is explicitly not protection against deliberate introspection;
an agent capable of arbitrary code execution needs a separate process with
limited inputs for an adversarial leakage test.

The root MIT notice and the vendored simglucose notice are present, and the
project provides attribution and upstream/vendored file hashes. Those records
make a future audit feasible; the inspected manifest is not an independent
verification of every upstream byte or a complete dependency and data-rights
review. The physiological mechanism has no claimed validity for depositors.
The immediate recommendation is to adopt the acquisition/evaluation contract
in the bank proposal, and evaluate any actual package reuse separately with
its pinned source, numerical environment and observation semantics.

\subsection{NeuronBench: explicit experiments with hidden operational fallbacks}

NeuronBench exposes deterministic and stochastic neuronal worlds, a protocol
menu and held-out intervention forecasts. Its documented distinction between
observed traces and hidden mechanism truth illustrates how an environment can
support active discovery without showing its answers to the learner. The
inspected source also includes a reference LLM baseline with an injected
client. That baseline is explicitly not Bayesian inference or numerical
value-of-information optimization
\citep{exec13_active_neuronbench}.

A source-level detail affects any comparison that reuses the baseline.
If an experiment response cannot be parsed, the implementation chooses the
first remaining protocol. If a forecast cannot be parsed, it returns the mean
observed count, or zero without observations. These behaviors keep the
procedure running but change the meaning of its output. A reported
``LLM-selected'' design may actually be a deterministic fallback. The proposed
adapter must therefore expose the attempted query, parse outcome, substituted
action, retry cost and downstream result. A useful ablation compares this
fallback-only policy with the model-assisted policy under equal environment
budgets. Errors should count in the assigned workload, while any
successful-only estimate receives its own explicit denominator.

The pinned project declares version 0.1.0, a Python requirement and an MIT
license. Its reference interface is a candidate for an optional public
simulator study; no package was installed or executed for this source
comparison. Porting its neurophysiological equations into the bank library
would add an unrelated model. Reusing an isolated environment interface, or
implementing the same interface for a transparent generated bank law, is the
more relevant proposal choice. Neither route establishes that the full
Model Discovery Agent has been integrated.

\subsection{DFM-Plans and unresolved implementation identities}

The inspected DFM-Plans repository describes itself as an overview, research
map and navigation point. Its README calls Zetema a research design and
states that individual implementations are released in separate repositories.
The pinned, nontruncated tree contains documentation and assets. Consequently,
this repository is unsuitable as a directly executable Zetema causal agent.
Its Apache-2.0 notice applies to the inspected repository; it does not confer
the rights or capabilities of separately linked projects
\citep{exec13_active_dfmplans}.

ScienceBuddy, ScienceIDE and JEPA-Anything are implementation leads from that
overview, not inspected capability evidence in this comparison. Likewise,
the benchmark repositories do not establish a complete public implementation
of the model-discovery method, and dataset links from the noisy-prior paper
do not establish its full agent package. A separately identified DISCERN
task and grading harness would also require source and dataset-rights review.
These distinctions prevent a literature survey from promising integration
with software that has not yet been identified at the required level.

The source map remains bounded by the stated literature cutoff and explicit
omissions. Later source inspection recovered the psychological
hypothesis-generation original, developed in
Section~\ref{exec20:hypothesis-scope}, and a scientific-method preprint,
discussed in Section~\ref{sec:exec24-scientific-method}. The latter does not
establish equivalence to the still-uninspected journal version. Additional
chemistry, hidden-structure and interactive-causality candidates remain
uninspected. No technical claim here relies on their discovery metadata, and
no absence or complete-frontier claim follows from those gaps. The original
scope record and the developed successor comparisons retain these distinct
coverage limits
\citep{exec13_active_scope}.

\section{A justified next experiment before an expensive bank trial}

The immediate empirical question is whether agent assistance improves the
translation from a decision brief to a valid acquisition plan. That can be
tested with generated worlds and public documentation. Begin with a brief
that specifies the decision, clock, admissible actions and loss, but withholds
the generator's mechanism. Compare a fixed design, a uniformly random design,
a numerical information criterion, a numerical decision-value criterion and
an optional agent-assisted proposal policy. The numerical policies need the
same candidate mechanisms and observations as the policy being compared; an
oracle with the true mechanism belongs in a separately labeled reference.

Vary the challenge along interpretable axes: irrelevant but highly informative
measurements, delayed outcomes, unavailable acquisitions, correlated expert
errors, observationally indistinguishable models, misspecified model pools
and malformed tool responses. Include clean cases that reward accepting
valid information. Evaluate the final decision loss, calibration of stated
probabilities where a probability target exists, the frequency of unsupported
causal claims, and whether constraints actually change the chosen action.
Graph accuracy alone cannot establish investment value.

Before final sampled runs, freeze the generated laws, held-out families,
permitted information, cost accounting, query and compute budgets, comparators,
failure handling and stopping rule. Report every assigned case and attempted
acquisition. If exploratory runs are used to select prompts or mechanisms,
keep them separate from the final comparison. A live model arm requires its
actual version, calls, prices and errors; an offline formulation exercise
supports claims about the design and reference logic only. The exact
construction above is already reproducible without live calls, random draws
or an expensive bank intervention.

A subsequent bank experiment becomes worth considering when the bounded
prototype identifies a consequential uncertainty that public and generated
evidence cannot settle. Its proposal should state which deployment decision
could change, the feasible observation law, the expected reduction in loss,
and the complete cost of obtaining that information. If candidate models
remain observationally equivalent under every permitted acquisition, the
appropriate result may be a restricted recommendation or abstention. An
agent adds value by recognizing that limit and specifying the missing
measurement or intervention, rather than by presenting another confident
causal narrative.
\label{exec13:active-experiments-end}


% SOURCE docs/release/2026-10-29/source/docs/writing-repair/2026-10-29/src/0f193208e9968c84.tex CHAR 0
% Preserving successor: the earlier chapter remains an inspected input.

